\documentclass[11pt,a4paper]{article}
% preamble.tex — shared setup for human-verification documents
% Compile each document with:  tectonic <File>.tex

\usepackage[T1]{fontenc}
\usepackage[utf8]{inputenc}
\usepackage{lmodern}
\usepackage{amssymb,amsmath,amsthm}
\usepackage{mathtools}
\usepackage[dvipsnames]{xcolor}
\usepackage{listings}
\usepackage{hyperref}
\usepackage[margin=2.5cm]{geometry}
\usepackage{parskip}

% ── Colours ──────────────────────────────────────────────────────────
\definecolor{leanBg}{HTML}{F7F7F0}
\definecolor{haskellBg}{HTML}{F0F4F7}
\definecolor{keyword}{HTML}{7B3294}
\definecolor{comment}{HTML}{5A8C32}
\definecolor{noteFg}{HTML}{B35900}

% ── Lean 4 listings language ─────────────────────────────────────────
% NOTE: Lean listing blocks use ASCII approximations of unicode symbols
% (e.g. -> for arrow, o. for compose).  The literate key substitutes
% these with proper math-mode renderings in the PDF.
\lstdefinelanguage{Lean4}{
  morekeywords={def,theorem,structure,where,namespace,end,import,fun,
                let,have,show,by,exact,intro,cases,rfl,subst,obtain,
                funext,simp,rw,apply,variable,Type,Prop,instance,
                section,open,class,extends,deriving,sorry,noncomputable},
  sensitive=true,
  morecomment=[l]{--},
  morecomment=[n]{/-}{-/},
  morestring=[b]",
}

% ── Lean listing style ───────────────────────────────────────────────
\lstdefinestyle{lean}{
  language=Lean4,
  backgroundcolor=\color{leanBg},
  basicstyle=\small\ttfamily,
  keywordstyle=\color{keyword}\bfseries,
  commentstyle=\color{comment}\itshape,
  breaklines=true,
  frame=single,
  framerule=0.4pt,
  xleftmargin=4pt,
  xrightmargin=4pt,
  aboveskip=6pt,
  belowskip=4pt,
  literate=
    {->}{{$\to$}}2
    {=>}{{$\Rightarrow$}}2
    {<|}{{$\langle$}}1
    {|>}{{$\rangle$}}1
    {==>}{{$\Longrightarrow$}}3
    { o. }{{$\;\circ\;$}}3
    {forall}{{$\forall$}}1
    {:=}{{$\,:=\,$}}2
    {!=}{{$\neq$}}2
    {<=}{{$\leq$}}2
    {>=}{{$\geq$}}2,
}

% ── Haskell listing style ───────────────────────────────────────────
\lstdefinestyle{haskell}{
  language=Haskell,
  backgroundcolor=\color{haskellBg},
  basicstyle=\small\ttfamily,
  keywordstyle=\color{keyword}\bfseries,
  commentstyle=\color{comment}\itshape,
  breaklines=true,
  frame=single,
  framerule=0.4pt,
  xleftmargin=4pt,
  xrightmargin=4pt,
  aboveskip=6pt,
  belowskip=4pt,
  literate=
    {->}{{$\to$}}2
    {=>}{{$\Rightarrow$}}2,
}

% ── Theorem environments ────────────────────────────────────────────
\theoremstyle{definition}
\newtheorem{defn}{Definition}[section]
\newtheorem{thm}[defn]{Theorem}
\newtheorem{prop}[defn]{Proposition}

\theoremstyle{remark}
\newtheorem*{fidelitynote}{Fidelity note}

% ── Convenience macros ───────────────────────────────────────────────
\newcommand{\lean}[1]{\texttt{#1}}      % inline Lean identifier
\newcommand{\hask}[1]{\texttt{#1}}      % inline Haskell identifier
\newcommand{\Ext}{\mathrm{Ext}}
\newcommand{\Pos}{\mathrm{Pos}}
\newcommand{\Shape}{\mathrm{Shape}}


\title{Human Verification: \texttt{Containers.Inventory}\\[4pt]
  \large Inventory as an Update Monad}
\author{Machine-generated from \texttt{lean/Inventory.lean} for review by Neil Ghani}
\date{July 2026}

\begin{document}
\maketitle

\begin{abstract}
This document presents every definition and theorem \emph{statement} from the
Lean~4 module \lean{Containers.Inventory} in three parallel forms: the
original Lean, a readable Haskell rendering, and precise mathematical English.
Proofs are omitted --- they are machine-checked by Lean.  The purpose is to
let a human verifier confirm that the formalisation captures the intended
mathematics.

The file verifies \S8 of the paper: directed containers with
\emph{constant position sets} are exactly \emph{monoid actions}, and these
co-interpret as \emph{update monads} $T X = S \to M \times X$.  Kleisli maps
observe a stock level in $S$ and emit a movement in $M$.  Two paper remarks
are machine-checked: abelianisation forgets audit-trail order, and monus
clamps on $\mathbb{N}$ while $\mathbb{Z}$ admits backorders.
\end{abstract}

\tableofcontents
\bigskip

%===================================================================
\section{Overview}
%===================================================================

In a directed container $\mathcal{D}$ where $\mathrm{Pos}(s) = M$ for all
shapes $s$ (positions independent of shape), the shift $\mathbf{shift}(s,
p, q)$ is genuinely a binary operation $M \times M \to M$ and does not depend
on $s$.  The five directed-container laws then collapse exactly to the axioms
of a monoid action: $M$ is the acting monoid, $S = \mathrm{Shape}$ is the set
of states.

The update monad $T_A X = A.S \to A.\mathrm{Mon}.M \times X$ is the
co-interpretation of this directed container: a Kleisli map looks at the
current state, emits a movement, and returns a value.


%===================================================================
\section{Monoid actions}
%===================================================================

% --- MonoidAction --------------------------------------------------
\subsection{Structure: \lean{MonoidAction}}

\begin{defn}[Bundled right monoid action]
A \textbf{bundled right monoid action} $A$ consists of:
\begin{itemize}
  \item a monoid $A.\lean{Mon}$ of \emph{stock movements};
  \item a type $A.S$ of \emph{states} (stock levels);
  \item an action $A.\lean{act} : A.S \to A.\lean{Mon}.M \to A.S$;
  \item a proof of \lean{act\_one}: $A.\lean{act}(s, 1) = s$; and
  \item a proof of \lean{act\_mul}: $A.\lean{act}(A.\lean{act}(s, m), m') =
    A.\lean{act}(s, m \cdot m')$.
\end{itemize}
\end{defn}

\paragraph{Lean.}
\begin{lstlisting}[style=lean]
structure MonoidAction where
  Mon     : Monoid'
  S       : Type
  act     : S -> Mon.M -> S
  act_one : forall s, act s Mon.one = s
  act_mul : forall s m m', act (act s m) m' = act s (Mon.mul m m')
\end{lstlisting}

\paragraph{Haskell.}
\begin{lstlisting}[style=haskell]
-- A right monoid action:
class Monoid m => Action s m where
  act :: s -> m -> s
  -- act_one: act s mempty = s
  -- act_mul: act (act s m) m' = act s (mappend m m')
\end{lstlisting}


%===================================================================
\section{Constant-position directed containers}
%===================================================================

% --- ConstDC -------------------------------------------------------
\subsection{Structure: \lean{ConstDC}}

\begin{defn}[Constant-position directed container]
A \textbf{constant-position directed container} $C$ has:
\begin{itemize}
  \item a type $C.S$ of shapes (states);
  \item a single position type $C.M$ (independent of the shape);
  \item $C.\lean{sub} : C.S \to C.M \to C.S$ (state transition);
  \item $C.\lean{root} : C.M$ (unit movement); and
  \item $C.\lean{shift} : C.M \to C.M \to C.M$ (movement composition).
\end{itemize}
Five laws (c1--c5) correspond to the directed-container laws D1--D5 with all
shape-dependent transports removed: the D2/D5 casts vanish because position
types are constant.
\end{defn}

\paragraph{Lean.}
\begin{lstlisting}[style=lean]
structure ConstDC where
  S     : Type
  M     : Type
  sub   : S -> M -> S
  root  : M
  shift : M -> M -> M
  c1 : forall s, sub s root = s
  c2 : forall m, shift root m = m
  c3 : forall m, shift m root = m
  c4 : forall s p q, sub s (shift p q) = sub (sub s p) q
  c5 : forall p q r, shift (shift p q) r = shift p (shift q r)
\end{lstlisting}

\paragraph{Haskell.}
\begin{lstlisting}[style=haskell]
data ConstDC s m = ConstDC
  { cdcSub   :: s -> m -> s
  , cdcRoot  :: m
  , cdcShift :: m -> m -> m
  -- c1..c5 as law comments
  }
\end{lstlisting}

\begin{fidelitynote}
The key insight of this section is that with constant positions, the shift
$C.\lean{shift} : C.M \to C.M \to C.M$ is a genuine binary operation on a
single set $M$ --- giving a monoid $(M, C.\lean{shift}, C.\lean{root})$ with
laws c2/c3/c5.  For varying $\mathrm{Pos}(s)$, the shift is only a fibrewise
composition and no global monoid exists.  This is the \emph{sound} case of
``the monoid that wasn't'' (noted in the dream journal).
\end{fidelitynote}


% --- ConstDC.toAction / MonoidAction.toConst ----------------------
\subsection{Definitions: \lean{ConstDC.toAction} and \lean{MonoidAction.toConst}}

\begin{defn}[Dictionary between \lean{ConstDC} and \lean{MonoidAction}]
The two structures are definitionally isomorphic:
\begin{itemize}
  \item $\lean{ConstDC.toAction}(C)$: the monoid is $(C.M, C.\lean{shift},
    C.\lean{root})$; the action is $C.\lean{sub}$; laws c2/c3/c5 give monoid
    laws; laws c1/c4 give action laws.
  \item $\lean{MonoidAction.toConst}(A)$: $S = A.S$, $M = A.\lean{Mon}.M$,
    $\lean{sub} = A.\lean{act}$, $\lean{root} = A.\lean{Mon}.\lean{one}$,
    $\lean{shift} = A.\lean{Mon}.\lean{mul}$.
\end{itemize}
\end{defn}

\paragraph{Lean.}
\begin{lstlisting}[style=lean]
def ConstDC.toAction (C : ConstDC) : MonoidAction where
  Mon := { M := C.M, mul := C.shift, one := C.root,
           mul_assoc := C.c5, one_mul := C.c2, mul_one := C.c3 }
  S := C.S; act := C.sub
  act_one := C.c1
  act_mul := fun s m m' => (C.c4 s m m').symm

def MonoidAction.toConst (A : MonoidAction) : ConstDC where
  S := A.S; M := A.Mon.M
  sub := A.act; root := A.Mon.one; shift := A.Mon.mul
  c1 := A.act_one; c2 := A.Mon.one_mul; c3 := A.Mon.mul_one
  c4 := fun s p q => (A.act_mul s p q).symm
  c5 := A.Mon.mul_assoc
\end{lstlisting}

\paragraph{Haskell.}
\begin{lstlisting}[style=haskell]
toAction :: ConstDC s m -> (s, m, m -> m -> m)
toAction c = (cdcSub c, cdcRoot c, cdcShift c)

toConst :: Action s m => s -> m -> ConstDC s m
toConst _ _ = ConstDC act mempty mappend
\end{lstlisting}


% --- round trips ---------------------------------------------------
\subsection{Theorems: \lean{MonoidAction.toConst\_toAction} and \lean{ConstDC.toAction\_toConst}}

\begin{thm}[Dictionary round trips]
The dictionaries are mutually inverse (definitionally):
\begin{enumerate}
  \item $A.\lean{toConst}.\lean{toAction} = A$.
  \item $C.\lean{toAction}.\lean{toConst} = C$.
\end{enumerate}
Both hold by \lean{rfl} (structure eta and proof irrelevance).
\end{thm}

\paragraph{Lean.}
\begin{lstlisting}[style=lean]
theorem MonoidAction.toConst_toAction (A : MonoidAction) :
    A.toConst.toAction = A := rfl

theorem ConstDC.toAction_toConst (C : ConstDC) :
    C.toAction.toConst = C := rfl
\end{lstlisting}

\paragraph{Haskell.}
\begin{lstlisting}[style=haskell]
-- toConst . toAction = id
-- toAction . toConst = id
\end{lstlisting}


% --- transport_const ----------------------------------------------
\subsection{Theorem: \lean{transport\_const}}

\begin{thm}[Constant-fibre transport is identity]
If the motive of a transport is constant (does not depend on the shape),
then the transport is the identity:
\[
  \forall\, h : s = t,\; q : M,\quad \mathrm{transport}_{M}(h, q) = q.
\]
This is the key lemma explaining why D2/D5 casts vanish in the constant-
position setting.
\end{thm}

\paragraph{Lean.}
\begin{lstlisting}[style=lean]
theorem transport_const {S M : Type} :
    forall {s t : S} (h : s = t) (q : M),
    @Eq.rec S s (fun _ _ => M) q t h = q
\end{lstlisting}

\paragraph{Haskell.}
\begin{lstlisting}[style=haskell]
-- In Haskell, transport along equalities does not exist
-- (no propositional equality at the term level).
-- This theorem is vacuously true: there is no Eq.rec to reason about.
\end{lstlisting}


% --- ConstDC.toDirected / MonoidAction.toDC -----------------------
\subsection{Definitions: \lean{ConstDC.toDirected} and \lean{MonoidAction.toDC}}

\begin{defn}[Embedding into DirectedContainer]
$\lean{ConstDC.toDirected}(C)$ embeds a constant-position directed container
into the general \lean{DirectedContainer} type by setting $\mathrm{Pos}(s) :=
C.M$ for all shapes $s$.  D1--D5 are exactly c1--c5 (with D2 and D5
augmented by \lean{transport\_const} to handle the casts).

$\lean{MonoidAction.toDC}(A) := A.\lean{toConst}.\lean{toDirected}$ is the
corresponding \lean{DirectedContainer} for a monoid action.
\end{defn}

\paragraph{Lean.}
\begin{lstlisting}[style=lean]
def ConstDC.toDirected (C : ConstDC) : DirectedContainer where
  Shape := C.S; Pos := fun _ => C.M
  root := fun _ => C.root
  sub := C.sub
  shift := fun _ p q => C.shift p q
  d1 := C.c1; d4 := C.c4; d3 := fun _ p => C.c3 p
  d2 := fun s q =>
    (C.c2 q).trans (transport_const (C.c1 s) q).symm
  d5 := fun s p q q' =>
    (congrArg (C.shift (C.shift p q))
              (transport_const (C.c4 s p q).symm q')).trans
      (C.c5 p q q')

def MonoidAction.toDC (A : MonoidAction) : DirectedContainer :=
  A.toConst.toDirected
\end{lstlisting}

\paragraph{Haskell.}
\begin{lstlisting}[style=haskell]
toDC :: Action s m => s -> DirectedContainer s m
toDC _ = DirectedContainer
  { dcPos   = const mempty  -- position type is always M
  , dcRoot  = mempty
  , dcSub   = act
  , dcShift = \_ -> mappend
  }
\end{lstlisting}


%===================================================================
\section{The update monad}
%===================================================================

% --- Update --------------------------------------------------------
\subsection{Definition: \lean{Update}}

\begin{defn}[Update monad]
The \textbf{update monad} of a monoid action $A$ is:
\[
  T_A\, X \;=\; A.S \to A.\lean{Mon}.M \times X.
\]
A value of type $T_A\, X$ is a \emph{Kleisli map}: it observes a current
state in $A.S$, emits a movement in $A.\lean{Mon}.M$, and returns a value
in $X$.  This is the co-interpretation of the constant-position directed
container $A.\lean{toDC}$ (Ahman--Uustalu, \emph{Update monads}).
\end{defn}

\paragraph{Lean.}
\begin{lstlisting}[style=lean]
def Update (A : MonoidAction) (X : Type) : Type :=
  A.S -> A.Mon.M * X
\end{lstlisting}

\paragraph{Haskell.}
\begin{lstlisting}[style=haskell]
newtype Update s m x = Update { runUpdate :: s -> (m, x) }
-- cf. Data.Functor.Product
\end{lstlisting}


% --- pure / bind ----------------------------------------------------
\subsection{Definitions: \lean{Update.pure} and \lean{Update.bind}}

\begin{defn}[Monad operations]
\begin{itemize}
  \item $\lean{pure}(A, x) : T_A\, X$ --- emit no movement, return $x$:
    $\lean{pure}(A, x)(s) = (1, x)$.
  \item $\lean{bind}(t, k) : T_A\, Y$ for $t : T_A\, X$, $k : X \to T_A\, Y$
    --- run $t$, apply the emitted movement to the state, run $k$ on the
    updated state, multiply the two movements:
    \[
      \lean{bind}(t, k)(s) \;=\;
      \bigl(t(s).\mathbf{1} \cdot k(t(s).\mathbf{2})(A.\lean{act}(s,\, t(s).\mathbf{1})).\mathbf{1},\;
      k(t(s).\mathbf{2})(A.\lean{act}(s,\, t(s).\mathbf{1})).\mathbf{2}\bigr).
    \]
\end{itemize}
\end{defn}

\paragraph{Lean.}
\begin{lstlisting}[style=lean]
def pure (A : MonoidAction) (x : X) : Update A X :=
  fun _ => (A.Mon.one, x)

def bind (t : Update A X) (k : X -> Update A Y) : Update A Y :=
  fun s =>
    (A.Mon.mul (t s).1 (k (t s).2 (A.act s (t s).1)).1,
     (k (t s).2 (A.act s (t s).1)).2)
\end{lstlisting}

\paragraph{Haskell.}
\begin{lstlisting}[style=haskell]
instance (Monoid m, Action s m) => Monad (Update s m) where
  return x = Update $ \_ -> (mempty, x)
  Update t >>= k = Update $ \s ->
    let (m,  x) = t s
        (m', y) = runUpdate (k x) (act s m)
    in  (mappend m m', y)
\end{lstlisting}


% --- Monad laws ----------------------------------------------------
\subsection{Theorems: monad laws}

\begin{thm}[Monad laws for \lean{Update}]
\begin{enumerate}
  \item \textbf{Left unit}: $\lean{bind}(\lean{pure}(A, x), k) = k(x)$.
    Uses \lean{act\_one} (the state is unchanged by the unit movement)
    and \lean{one\_mul}.
  \item \textbf{Right unit}: $\lean{bind}(t, \lean{pure}(A)) = t$.
    Uses \lean{mul\_one}.
  \item \textbf{Associativity}: $\lean{bind}(\lean{bind}(t, k), h) =
    \lean{bind}(t, \lambda x.\, \lean{bind}(k(x), h))$.
    Uses \lean{act\_mul} and \lean{mul\_assoc}.
\end{enumerate}
\end{thm}

\paragraph{Lean.}
\begin{lstlisting}[style=lean]
theorem pure_bind (x : X) (k : X -> Update A Y) :
    bind (pure A x) k = k x

theorem bind_pure (t : Update A X) : bind t (pure A) = t

theorem bind_assoc (t : Update A X) (k : X -> Update A Y)
    (h : Y -> Update A Z) :
    bind (bind t k) h = bind t (fun x => bind (k x) h)
\end{lstlisting}

\paragraph{Haskell.}
\begin{lstlisting}[style=haskell]
-- Standard monad laws:
-- return x >>= k      = k x         (left unit)
-- t >>= return        = t            (right unit)
-- (t >>= k) >>= h     = t >>= (\x -> k x >>= h)   (associativity)
\end{lstlisting}


%===================================================================
\section{Paper remarks, machine-checked}
%===================================================================

% --- abelianization -----------------------------------------------
\subsection{Example: abelianisation forgets audit-trail order}

\begin{prop}[Abelianisation forgets order]
Two movement words with the same net balance can differ as audit trails, and
can have different operational consequences.  Specifically:
\[
  \mathrm{net}[\mathrm{receipt}(2), \mathrm{issue}(1)]
  \;=\; \mathrm{net}[\mathrm{issue}(1), \mathrm{receipt}(2)]
\]
(both equal 1), yet $[\mathrm{receipt}(2), \mathrm{issue}(1)] \neq
[\mathrm{issue}(1), \mathrm{receipt}(2)]$ as lists.  The second word passes
through a shortage (from zero stock) that the first never sees.
\end{prop}

\paragraph{Lean.}
\begin{lstlisting}[style=lean]
example :
    net [.receipt 2, .issue 1] = net [.issue 1, .receipt 2] /\
    ([Movement.receipt 2, .issue 1] != [.issue 1, .receipt 2]) := by
  decide
\end{lstlisting}

\paragraph{Haskell.}
\begin{lstlisting}[style=haskell]
-- net [Receipt 2, Issue 1] = net [Issue 1, Receipt 2]  -- both = 1
-- but [Receipt 2, Issue 1] /= [Issue 1, Receipt 2]     -- as lists
\end{lstlisting}


% --- monus clamps / Z admits backorders ---------------------------
\subsection{Theorems: monus clamps; $\mathbb{Z}$ admits backorders}

\begin{prop}[Monus clamps on $\mathbb{N}$]
Truncated subtraction on $\mathbb{N}$ silently discards unmet demand:
$\exists s, m : \mathbb{N},\; (s - m) + m \neq s$.
Witness: $s = 0$, $m = 1$.
\end{prop}

\begin{thm}[$\mathbb{Z}$ admits backorders]
For all $s, m : \mathbb{Z}$, $(s - m) + m = s$.
\end{thm}

\paragraph{Lean.}
\begin{lstlisting}[style=lean]
example : exists s m : Nat, (s - m) + m != s := <|0, 1, by decide|>

theorem int_sub_add_cancel : forall s m : Int, (s - m) + m = s := by
  intro s m; omega
\end{lstlisting}

\paragraph{Haskell.}
\begin{lstlisting}[style=haskell]
-- Nat monus: (0 - 1) + 1 = 0 /= 0+1 = 1  -- no, (0 - 1) = 0 in Nat
-- so  (0 - 1) + 1 = 1 /= 0   -- True: Nat monus clamps
monusClamps :: Bool
monusClamps = ((0 :: Word) - 1) + 1 /= 0  -- True

-- Int: (s - m) + m = s always
intCancel :: Integer -> Integer -> Bool
intCancel s m = (s - m) + m == s
\end{lstlisting}


%===================================================================
\section{Concrete instance: two-product warehouse}
%===================================================================

% --- intPairMonoid / inventory -----------------------------------
\subsection{Definitions: \lean{intPairMonoid} and \lean{inventory}}

\begin{defn}[Integer-pair monoid and inventory action]
The \textbf{integer-pair monoid} is $(\mathbb{Z} \times \mathbb{Z}, +, (0,0))$
(componentwise addition).  The \textbf{inventory action} is the right action
of $\mathbb{Z}^2$ on $\mathbb{Z}^2$ by componentwise addition:
$\lean{act}((s_1, s_2), (m_1, m_2)) = (s_1 + m_1, s_2 + m_2)$.
Negative stock levels represent backorders (accepted-but-unfulfilled demand).
\end{defn}

\paragraph{Lean.}
\begin{lstlisting}[style=lean]
@[reducible] def intPairMonoid : Monoid' where
  M := Int * Int
  mul a b := (a.1 + b.1, a.2 + b.2)
  one := (0, 0)
  ...

@[reducible] def inventory : MonoidAction where
  Mon := intPairMonoid
  S := Int * Int
  act s m := (s.1 + m.1, s.2 + m.2)
  ...
\end{lstlisting}

\paragraph{Haskell.}
\begin{lstlisting}[style=haskell]
type Stock = (Integer, Integer)

instance Semigroup Stock where
  (a1, a2) <> (b1, b2) = (a1+b1, a2+b2)

instance Monoid Stock where
  mempty = (0, 0)

instance Action Stock Stock where
  act (s1, s2) (m1, m2) = (s1+m1, s2+m2)
\end{lstlisting}


% --- warehouseDC / policy -----------------------------------------
\subsection{Definitions: \lean{warehouseDC}, \lean{receive}, and \lean{policy}}

\begin{defn}[Warehouse DC, replenishment policy]
The \textbf{warehouse directed container} \lean{warehouseDC} $:=
\lean{inventory}.\lean{toDC}$ is the directed container with stock levels as
shapes and $\mathbb{Z}^2$-movements as positions.

\lean{receive}$(m) : T_{\lean{inventory}}\, ()$ is a Kleisli map that
unconditionally emits the movement $m$.

The \textbf{order-up-to replenishment policy} is a Kleisli map: for each
product, if stock is below~5, order up to~10:
\[
  \lean{policy}((), s) \;=\; \bigl((\text{if } s_1 < 5 \text{ then } 10 - s_1
  \text{ else } 0,\; \text{if } s_2 < 5 \text{ then } 10 - s_2 \text{ else } 0),\;
  ()\bigr).
\]
\end{defn}

\paragraph{Lean.}
\begin{lstlisting}[style=lean]
def warehouseDC : DirectedContainer := inventory.toDC

@[reducible] def receive (m : Int * Int) : Update inventory Unit :=
  fun _ => (m, ())

@[reducible] def policy : Unit -> Update inventory Unit :=
  fun _ s =>
    ((if s.1 < 5 then 10 - s.1 else 0,
      if s.2 < 5 then 10 - s.2 else 0), ())
\end{lstlisting}

\paragraph{Haskell.}
\begin{lstlisting}[style=haskell]
type WarehouseM = Update Stock Stock

receive :: Stock -> WarehouseM ()
receive m = Update $ \_ -> (m, ())

policy :: () -> WarehouseM ()
policy () = Update $ \(s1, s2) ->
  ((if s1 < 5 then 10 - s1 else 0,
    if s2 < 5 then 10 - s2 else 0), ())
\end{lstlisting}


\subsection{Sanity checks (examples checked by \lean{decide})}

\begin{prop}[Policy with product 1 low, product 2 high]
With stocks $(2, 7)$ (product~1 below reorder, product~2 above):
\begin{itemize}
  \item the policy emits $(8, 0)$ (order~8 of product~1, nothing of~2);
  \item applying the emitted order brings stocks to $(10, 7)$.
\end{itemize}
\end{prop}

\begin{prop}[Kleisli composition threads state]
Receiving 3 units of product~1 before the policy, starting from stocks
$(1, 7)$: the policy sees stock~4 (= $1 + 3$) and orders~6; total emitted
movement is $(9, 0)$.
\end{prop}

\paragraph{Lean.}
\begin{lstlisting}[style=lean]
example : (policy () (2, 7)).1 = (8, 0) := by decide
example : inventory.act (2, 7) (policy () (2, 7)).1 = (10, 7) := by decide
example :
    Update.bind (receive (3, 0)) (fun _ => policy ()) (1, 7)
      = ((9, 0), ()) := by decide
\end{lstlisting}

\paragraph{Haskell.}
\begin{lstlisting}[style=haskell]
-- runUpdate (policy ()) (2, 7) = ((8, 0), ())
-- act (2,7) (8,0) = (10,7)
-- runUpdate (receive (3,0) >>= (\_ -> policy ())) (1,7) = ((9,0),())
\end{lstlisting}


%===================================================================
\section{Translation conventions}
%===================================================================

\begin{enumerate}
\item \textbf{Bundled vs.\ typeclass.}  Lean uses bundled records for
  \lean{Monoid'} and \lean{MonoidAction}; Haskell uses typeclasses.  The
  mathematical content is the same.

\item \textbf{Transport in D2/D5.}  In the general directed-container
  formalism, D2 and D5 involve transports along shape equalities.  In the
  constant-position case, these transports are the identity (\lean{transport\_const}),
  so c2 and c5 are the bare monoid laws without casts.

\item \textbf{Backorders as negative stock.}  The paper notes that $\mathbb{N}$
  monus clamps while $\mathbb{Z}$ is honest.  The Lean formalisation uses
  $\mathbb{Z}$ throughout the warehouse instance.

\item \textbf{Proof bodies omitted.}  Only definitions and theorem
  \emph{statements} are translated.  Lean checks the proofs; the human
  verifier checks that the \emph{statements} match the intended mathematics.
\end{enumerate}

\end{document}
